\documentclass[../main.tex]{subfiles}

\begin{document}

\chapter*{\myAbstractTitle}

This thesis investigates whether political party positions in election manifestos predict technology diffusion across democracies, addressing three research questions: how to reliably classify political positions, what determines technology diffusion, and whether political positions impact diffusion. \par

The primary analysis employs panel data regression combining the Cross-Country Historical Adoption of Technology (CHAT) dataset with Comparative Manifestos Project (CMP) coding, concluding in over 100 technologies and 34 democracies in the timeframe between 1950 and 2000. 
CMP variables are aggregated into "literature aspects" linking political rhetoric to diffusion determinants. 
The panel model incorporates country and technology-year fixed effects. 
A robustness check employs zero-shot Large Language Model (LLM) classification for migration policy. \par

The central finding is that manifesto positions do not robustly predict technology diffusion. 
After False Discovery Rate correction, no aggregated literature aspects show significant relationships across sectors or time delays (0-4 years). 
This null result holds across both CMP-based and independently validated LLM-generated measures. 
Panel regressions exhibit high adjusted R$^2$ (above 80\%) but low within-R$^2$ values (averaging 2\%), revealing that stable country characteristics and global technology trends explain most variation, while time-varying political rhetoric explains virtually none of the within-country variation. \par

This pattern suggests politics influences diffusion through stable institutional frameworks (e.g., property rights, education) that change slowly over decades and are captured by fixed effects, rather than manifesto variations. 
The null finding reflects temporal mismatch between diffusion and electoral cycles, implementation gaps, and structural constraints limiting rapid changes. \par

Methodologically, zero-shot LLM classification proves viable for political text coding. 
Gemini 2.5 Pro achieved 47.2\% accuracy replicating CMP codes (competitive with fine-tuned models) and 0.88 correlation with expert surveys for migration (outperforming traditional CMP and other approaches). 
Both measurement approaches converged on the same conclusion. \par

Electoral manifestos capture political rhetoric but fail to predict short-term technology diffusion. 
Politics shapes adoption through stable institutional frameworks and continuous engagement, not electoral rhetoric.

\vspace{1cm}
\textbf{Keywords:} Technology Diffusion, Political Manifestos, Large Language Models, Zero-Shot Classification, Panel Data Analysis, Comparative Manifestos Project, CHAT Dataset

\end{document}